\section{Latent Diffusion \& Diffusion Transformers (DiT)}\label{sec:ldm}
\lab{LDM \textperiodcentered\ Stable Diffusion \textperiodcentered\ VAE compression \textperiodcentered\ DiT \textperiodcentered\ adaLN-Zero \textperiodcentered\ SDXL}

\noindent\begin{minipage}[t]{0.42\columnwidth}\vspace{0pt}\centering
\resizebox{\linewidth}{!}{%
\begin{tikzpicture}[node distance=1.1mm,every node/.style={font=\tiny},
  bx/.style={box,minimum width=21mm,minimum height=3.3mm,inner sep=1pt}]
\node[bx,fill=soft] (img) {Pixel Image $512\times 512\times 3$};
\node[bx,draw=acc,above=1.5mm of img] (vae) {Encoder $\mathcal{E}$ ($f{=}8$ downsample)};
\node[bx,draw=acc2,fill=soft2,above=1.5mm of vae] (lat) {Latent $z \in \mathbb{R}^{64\times 64\times 4}$};
\node[bx,draw=acc3,fill=acc3!8,above=1.5mm of lat] (dit) {Diffusion UNet / DiT};
\node[bx,above=1.5mm of dit] (dec) {Decoder $\mathcal{D}$ (Latent $\to$ Pixels)};
\draw[arr] (img)--(vae); \draw[arr] (vae)--(lat); \draw[arr] (lat)--(dit); \draw[arr] (dit)--(dec);
\end{tikzpicture}%
}
\end{minipage}\hfill
\begin{minipage}[t]{0.56\columnwidth}\vspace{1pt}\small
\textbf{Perceptual vs Semantic Compression.} Pixel space diffusion wastes $>90\%$ of capacity on imperceptible high-frequency pixel noise.\\[2pt]
{\color{acc2}$\blacktriangleright$} \textbf{$f=8$ Compression}: $512\times 512\times 3 = 786{,}432$ values compressed to $64\times 64\times 4 = 16{,}384$ values ($48\times$ memory reduction).\\
{\color{acc2}$\blacktriangleright$} \textbf{Cross-modal injection}: Text conditioning injected into spatial latents via cross-attention.
\end{minipage}

\begin{mem}[Latent Diffusion Models Compared]
\begin{tabularx}{\linewidth}{@{}p{16mm} c p{22mm} X@{}}
\toprule
\textbf{Model} & \textbf{Params} & \textbf{Text Encoders} & \textbf{Conditioning / Backbone}\\
\midrule
\textbf{SD 1.5} & 860M & CLIP ViT-L/14 & UNet $512^2$, Cross-attention\\
\textbf{SD 2.1} & 865M & OpenCLIP-ViT-H & UNet $768^2$, $v$-pred, cross-attn\\
\textbf{SDXL}   & 2.6B & CLIP-L + OpenCLIP-G & UNet $1024^2$, Dual text, crop coords\\
\textbf{DiT-XL/2} & 675M & Class / CLIP tokens & Transformer $256^2$, adaLN-Zero\\
\bottomrule
\end{tabularx}
\end{mem}

\begin{qa}[Architecture Shifts: UNet to DiT]
\Q{Why did modern generative vision abandon convolutional UNets for Diffusion Transformers (DiT)?}{
\textbf{1. Scaling Laws}: Convolutional UNets plateau in generation quality as parameter counts scale past 2B due to fixed receptive fields and localized inductive biases. Transformers scale monotonically with compute and FLOPs without saturation.
\textbf{2. Flexible sequence handling}: UNets require rigid spatial grid tensors. DiT patchifies latents into 1D token sequences, making it trivial to process arbitrary aspect ratios, variable image resolutions, and temporal video frames in a single unified architecture.
\textbf{3. Hardware efficiency}: Standard dense matrix multiplications in Transformer self-attention achieve $>50\%$ MFU on Tensor Cores (especially with FlashAttention-2), whereas multi-scale UNet up/down convolutions suffer from high memory bandwidth overhead.}
\Q{How does adaLN-Zero work in DiT, and why is zero-initialization necessary?}{
Rather than injecting timestep and class embeddings via cross-attention, DiT uses \textbf{Adaptive Layer Normalization (adaLN)}:
A projection MLP maps the sum of timestep and condition embeddings $(t + c)$ into scale, shift, and gate parameters $(\gamma_1, \beta_1, \alpha_1, \gamma_2, \beta_2, \alpha_2)$:
\begin{align*}
h' &= h + \alpha_1 \cdot \text{Attention}\big((1 + \gamma_1)\odot \text{LN}(h) + \beta_1\big)\\
h_{\text{out}} &= h' + \alpha_2 \cdot \text{MLP}\big((1 + \gamma_2)\odot \text{LN}(h') + \beta_2\big)
\end{align*}
\textbf{Zero-Initialization}: The gating factors $\alpha_1, \alpha_2$ are initialized to \textbf{zero}.
At initialization, every DiT block acts as an exact identity function ($h_{\text{out}} = h$). Gradients flow cleanly through arbitrarily deep networks at step 0 without exploding or vanishing.}
\Q{How does SDXL eliminate training crop artifacts via micro-conditioning?}{
Previous diffusion models discarded non-square images by center-cropping them to $512\times 512$, teaching the model that subjects are often cut in half at image borders.
\textbf{SDXL Micro-Conditioning (Podell et al., 2023)}:
During training, random crops are drawn from variable aspect-ratio images, and the exact pixel crop coordinates $(c_{\text{top}}, c_{\text{left}})$ along with original resolution $(h_{\text{orig}}, w_{\text{orig}})$ are embedded via Fourier features and injected into the UNet conditioning vector.
At inference, setting $(c_{\text{top}}=0, c_{\text{left}}=0)$ conditions the model to generate fully uncropped, centered compositions.}
\end{qa}

\begin{trap}[Latent Diffusion Pitfalls]
\begin{itemize}
\item \textbf{VAE Latent Scaling Factor}: Latent representations from the VAE encoder have non-unit variance ($\sigma \ne 1$). SD scales latents by constant factor $0.18215$ (SDXL uses $0.13025$) before adding noise: $z = 0.18215 \cdot \mathcal{E}(x)$. Forgetting this scaling factor leads to complete generation failure (snow noise).
\item \textbf{Latent VAE artifacts}: High spatial compression ($f=8$) creates artifacts on fine typography and small faces. Face restorers (CodeFormer) or lower compression ($f=4$) are required for micro-detail reconstruction.
\end{itemize}
\end{trap}
